\documentclass[12pt, a4paper]{article}

\usepackage{amssymb}
\usepackage{amsfonts}
\usepackage{amsmath}
\usepackage{amsthm}

\usepackage{mathtools}
\usepackage{cancel}
\usepackage[T1]{fontenc}
\usepackage[document]{ragged2e}
\usepackage{fancyhdr} % for header and footer
\usepackage{tcolorbox}
%\usepackage{graphicx} %for inserting picture
%\usepackage{float} % eh centering?
%\usepackage{subfigure} % for side by side figure
%\usepackage{array} % for fixed table
\usepackage{palatino}
\usepackage{titlesec}
\usepackage{indentfirst}

\usepackage[a4paper, top=3cm, bottom=3.5cm, left=3cm, right=3cm]{geometry}

\pagestyle{plain}

\setlength{\parindent}{1.5em} 
\linespread{1.15}

\newtheorem{theorem}{Theorem}[section]
\newtheorem{lemma}[theorem]{Lemma}
\theoremstyle{definition}
\newtheorem{definition}[theorem]{Definition}
\newtheorem*{example}{Example}

% Mengubah text default "Proof" menjadi "Bukti"
%\addto\captionsindonesian{\renewcommand{\proofname}{Bukti}}

\titleformat{\section}[block]
  {\normalfont\normalsize\bfseries\centering\scshape} 
  {\thesection.}                                     
  {0.5em}
  {}
\titlespacing*{\section}{0pt}{3em}{1.5em}     

\titleformat{\subsection}[block]
  {\normalfont\normalsize\bfseries\scshape}
  {\thesubsection}
  {0.5em}
  {}
\titlespacing*{\subsection}{0pt}{2em}{1em}

\nopagebreak
\fancypagestyle{plain}{
  \fancyhf{}
  \fancyfoot[C]{\thepage}                          
  \renewcommand{\headrulewidth}{0pt}
}

\fancypagestyle{firstpage}{
  \fancyhf{}
  \fancyhead[L]{
    \small\scshape Miles Reid Notes\\
    \small Undergraduate Algebraic Geometry \\
    \small FMIPA ITB
  }
  \fancyfoot[C]{\thepage}
  \renewcommand{\headrulewidth}{0pt}
}

\begin{document}
    \thispagestyle{firstpage}

    \begin{center}
        \vspace*{0.8cm}
        {\large\bfseries Miles Reid NOTE} \\[0.6cm]
        {\normalsize for and by Gilbert Allister (10123021)}
        \vspace*{0.8cm}
    \end{center}



    \section{Motivation}
    This one is reference to Miles Reid Notes, for me to easily read and memorized in my mind. This one also summarization to Week 1 until Week 5 of my lecturer's consultation. 
    
    \section{Affine Varieties}
    \subsection{Noetherian Rings}
    \begin{definition}[Noetherian]
        A ring $R$ is called Noetherian if every ideal of $R$ is finitely generated.     
    \end{definition}

    \begin{theorem}[Hilbert Basis Theorem]
        Let $R$ be a Noetherian ring. Then the polynomial ring $R[x]$ is also Noetherian.
    \end{theorem}


    \section{Zariski Topology}
    \subsection{Algebraic Set}

    \subsection{The Corresponding of Algebraic Set}

    \subsection{Irreducible Algebraic Set}


    \section{Nullstellensatz}
    \subsection{Radical Ideals}
    \subsection{Nullstellensatz theorems}
    Or also know as Hilbert's zeros theorem. 
    \begin{theorem}[Weak Nullstellensatz, geometric form]
      Suppose $k$ is an algebraically closed field. Let $I$ be a proper ideal in the polynomial ring $k[x_1, \ldots, x_n]$. Then we have 
      \begin{align*}
        V(I) \neq \emptyset
      \end{align*} 
      i.e., there exists a point $(a_1, \ldots, a_n) \in k^n$ such that $f(a_1, \ldots, a_n) = 0$ for all $f \in I$.
    \end{theorem}


    \begin{theorem}[Weak Nullstellensatz, algebraic form]
      Suppose $k$ is an algebraically closed field, and let $\mathfrak{m}$ is maximal ideal of $k[x_1, \ldots, x_n]$, then we have
      \begin{align*}
        \mathfrak{m} = (x_1-a_1,\ldots, x_n-a_n)
      \end{align*}
    \end{theorem}

    \begin{theorem}[Strong Nullstellensatz]
      Suppose $k$ is an algebraically closed field. Let $I$ be an ideal in the polynomial ring $k[x_1, \ldots, x_n]$. Then the ideal of all polynomials that vanish on the variety defined by $I$ is equal to the radical of $I$, i.e.,
      \begin{align*}
        I(V(I)) = \sqrt{I}.
      \end{align*}
    \end{theorem}

    \subsection{Corollary of Nullstellensatz}


\end{document}


